In this section we want to address a bit more formally two closely related questions:
\begin{enumerate}
	\item What is the exact space of solutions for the solutions to the basic model? In the definition \ref{ch2:def:solution-model} did not specify the codomain of the \rvs involved.
	\item How can we be sure that the conditioning $\Prb$ we are obtaining \emph{regular} conditional probability measures (\cite[Chapter 10.2,first definition]{dudley}), \ie that we obtain mathematical objects that behave as probability measures?
	This is an important question because for the definition \ref{ch2:def:solution-model} we need certain objects to be \Bms under $\PrbCondTimes{\cdot}{}$, thus we assume that $\PrbCondTimes{\cdot}{}$ is a probability measure.
\end{enumerate}

We will not give a complete answer to the second question, but we will give a direction to search for it, and we expect that with further work one might prove the necessary technicalities required to obtain the desired result.

\subsection{The space of solutions}\label{app:formal-def-solution}

We first need to check where the spawning Poisson processes $(\hmactive_t^p)_p$ are taking values for each fixed time $t$.
This is of course $(\N_0)^N$.
Now we define the space of configurations for the cAMP molecules $(\camp^{p,g})_g$ spawned by a single amoeba $\amba^p$.
Here the symbol "$\spadesuit$" means that a particle has not been spawned yet.
We must take only sequences where only the first coordinates (finitely) are different from $\spadesuit$.

\begin{definition}
Define the space of the positions cAMP particles spawned by an amoeba as:
\begin{equation*}
\confCamp(d):=\boxplus_{g=1}^\N (\{\spadesuit\} \sqcup \R^d) := 
\{
	x = (x_g)_{g\in\N}\
	:\
	\exists G\in\N_0.\
	\forall g < G.\
	x_g \in \R^d\
	\text{ and }\
	\forall g\geq G.\
	x_g =\spadesuit
\}
\end{equation*}
Define the space of configurations for all particles (the $N$ amoebas $\amba^p$ and the cAMP molecules $\camp^{p,g}$) as
\begin{equation*}
\Conf:=(\R^d)^N \times\big(\confCamp(d)\big)^N
\end{equation*}
For each configuration of particles $x\in \Conf$ sign its $d$-dimensional coordinates as 
\begin{equation*}
\begin{aligned}
x &= (x^1, \ldots, x^N, 
	&(x^{1,g})_{g\in\N}, \ldots, 
	&(x^{p,g})_{g\in\N}, \ldots, 
	&(x^{N,g})_{g\in\N})
\\
&= (a^1, \ldots, a^N, 
	&c^{1,1}, c^{1,2}, \ldots, 
	&c^{p,1},\ldots, 
	&c^{N,1}, \ldots)
\end{aligned}
\end{equation*}
where the absence of a coordinate in the list means that such coordinate is set to $\spadesuit$.
\end{definition}

\begin{remark}
	Notice that, if $\{\spadesuit\} \sqcup \R^d$ could be seen as an algebraic structure with $\spadesuit$ as a neutral element, then
	\begin{align*}
		\confCamp(d)
		&:=\boxplus_{g=1}^\N (\{\spadesuit\} \sqcup \R^d)\
		\subseteq
		\\
		&\subseteq\
		\bigoplus_{g=1}^\N (\{\spadesuit\} \sqcup \R^d)\
		:=
		\{
			x = (x_g)_{g\in\N}\
			:\
			\exists G\in\N_0.\
			\forall g\geq G.\
			x_g =\spadesuit
		\}
	\end{align*}
	And the latter one, among the other things, would be a categorical construct on which one can put quite a lot of structure (categorical product of an in algebraic category), which would be useful in case we would like to inquire about properties of such spaces.
	It does not seem to us though that there is any "natural" way to put an algebraic structure on $\{\spadesuit\} \sqcup \R^d$.
\end{remark}

And what about the Brownian motion associated to each particle? That is just a countable amount of independent $d$-dimensional \Bms.
With these definitions the solution of \ref{ch2:def:solution-model} can be seen as a \emph{random \cadlag} function 

\begin{equation*}
	((\hmactive^p)_p, \sol = (\amba, \camp), B) (\omega):\
	\R_+ \longrightarrow 
	\underbrace{
		(\N_0)^N
	}_{\text{Poisson processes}}
	\times
	\underbrace{
		(\R^d\times\confCamp(d))^N 
	}_{ \text{positions particles}}
		\times 
	\underbrace{
		(\R^d)^N \times (\R^d)^\N
	}_{\Bms}
\end{equation*}

or, in other words, the \rv $((\hmactive^p)_p, \sol, B)$ takes values in the set

\begin{equation}\label{app:def:space-of-solutions}
	\left\{
		f:\R_+ \longrightarrow 
			(\N_0)^N
			\times
			(\R^d\times\confCamp(d))^N
			\times 
			(\R^d)^N \times (\R^d)^\N\
			:\
			f \text{ is \cadlag}
	\right\}
\end{equation}

It is not necessary to redefine the drift $\drift$ and the diffusion $\diffusion$, because they become relevant for the SDE framework only under the conditioning $\PrbCondTimes{}{}$. In any case, if we really wanted to see them in this framework, we would have for absent particle they have no mapping, so one could set them to  $\spadesuit$ for every cAMP molecule not yet spawned and obtain consistent mappings in the appropriate spaces.

\subsection{Can we obtain regular probability measures by conditioning $\Prb$?}

The main instrument to verify that a measure is regular is \cite[Theorem 10.2.2]{dudley}, namely:

\begin{theorem}
	Let $T$ be any Polish space, $\Borel$ its $\sigma$-algebra of Borel sets.
	Let $(\Omega, \sigAlg, \Prb)$ be a probability space, $Y:\Omega\longrightarrow T$ a $(\sigAlg,\Borel)$-measurable function and $\mathcal{C}$ any sub-$\sigma$-algebra of $\sigAlg$.
	Then there exists a conditional distribution $\Prb_{Y|\mathcal{C}}(\cdot, \cdot \cdot):\Borel \times \Omega \longrightarrow [0,1]$ (see \cite[Chapter 10.2,first definition]{dudley}) and it is unique, meaning that if another $\Prb'(\cdot, \cdot \cdot)$ satisfies the definition of conditional distribution, then
	\begin{equation*}
		\forall^\Prb \omega\in\Omega.
		\quad
		\Prb_{Y|\mathcal{C}}(\cdot, \omega) = \Prb'(\cdot, \omega)
	\end{equation*}
\end{theorem}

[[correct]]

[[use theorem 2.8.2 from dudley to add the point to Rd and use the paragraph 13.2 to have that the product of such spaces can be seen as a polish space.
use chapter 2.4 from dudley to put a metric on the function spaces for the hmactive processes, if necessary.
Or should I maybe see this as a putting a function from R+ to the configuration spaces and N times the natural numbers (for the thing).
Maybe use theorem 2.4.5 from dudley and see, for every T>0 the space [0,T] as the compact space, and then the rest of the space a metric maybe polish space
You can actually see it as a skochorod space (which is polish) over a polish space, and thus a polish space. Complicated but doable]]